% !PRINTED 213
\weberpage{213}
Um die übrigen Substitutionen der Gruppe zu finden, bemerken
 wir, dass wir nach dem Theorem \S.\,53, 2., wenn wir $\pm 1$
für $a, a'$ und $0, \infty$ für $b, b'$ nehmen, auf die Existenz einer Substitution
 $\chi$ in $G$ schliessen können, die den Bedingungen genügt
\[
\chi(-1) = 0, \quad \chi(+1) = \infty,
\]
und dass also $\chi$ die Form haben muss:
\weq{2}{\chi = \lambda\,\frac{x+1}{x-1},}
wo $\lambda$ ein constanter Factor ist.

Wenn wir in $\chi$ die sechs conjugirten Pole $0, \infty, +1, -1,$
$+i, -i$ einsetzen, so müssen wir dieselben Werthe in einer
anderen Reihenfolge erhalten (\S.\,52, 5.). Diese Werthe sind aber
\[
-\lambda,\ \lambda,\ \infty,\ 0,\ -\lambda i,\ \lambda i.
\]

Es muss also $\lambda = \pm 1$ oder $= \pm i$ sein. Der Werth $\pm 1$
ist nicht zulässig, weil sonst
\[
\chi(i) = \mp i, \quad \Theta\chi(i) = \pm i,
\]
und also $\pm i$ die Pole von $\chi$ oder von $\Theta\chi$ wären, während sie
doch nur zweizählig und die Pole von $\Theta\psi$ sind. Es muss also
$\lambda = \pm i$ sein, und wir können ohne Beschränkung der Allgemeinheit
 das obere Zeichen nehmen, denn das untere Zeichen
entspricht dann der Substitution $\Theta\chi$. Es ist daher
\[
\chi(x) = i\,\frac{x+1}{x-1}.
\]

Daraus erhält man
\[
\chi^2(x) = \frac{x+i}{x-i}, \quad \chi^3(x) = x,
\]
und es ergeben sich die 12 Substitutionen:
\weq{3}{\begin{array}{cccc}
1, & \Theta, & \psi, & \Theta\psi \\
\chi, & \Theta\chi, & \psi\chi, & \Theta\psi\chi \\
\chi^2, & \Theta\chi^2, & \psi\chi^2, & \Theta\psi\chi^2,
\end{array}}
oder
\weq{4}{\pm x,\ \pm\frac{1}{x},\ \pm i\,\frac{x+1}{x-1},\ \pm i\,\frac{x-1}{x+1},\ \pm\frac{x+i}{x-i},\ \pm\frac{x-i}{x+i}.}

Dass diese Substitutionen wirklich eine Gruppe, die \gesperrt{Tetraëdergruppe},
bilden, von der $\Theta, \psi, \chi$ die erzeugenden Elemente
sind, ergiebt sich aus
\weq{5}{\chi\Theta = \Theta\psi\chi,\ \chi^2\Theta = \psi\chi^2,\ \psi\Theta = \Theta\psi,\ \chi\psi = \Theta\chi,\ \chi^2\psi = \Theta\psi\chi^2,}
%CONT>>
